\begin{abcliste}
\abc The impedance of a parallel resonant circuit has its peak at the resonance frequency. The peak lies at $f_0 \approx \result{\fzO}$: the radio receives the station on this frequency.

\abc At the resonance frequency,
\begin{align}
 2\pi f_0 = \frac{1}{\sqrt{L C}} \quad\Rightarrow\quad C &= \CF \\
   &= \frac{1}{\left(2\pi\times\fz\right)^2\times\Lr} \\
   &= \C \approx \result{\CP{-12}{2}}
\end{align}

\abc The same coil with another frequency:
\begin{align}
 C_2 &= \CbF \\
   &= \frac{1}{\left(2\pi\times\fb\right)^2\times\Lr} \\
   &= \Cb \approx \result{\CbP{-12}{2}}
\end{align}
From \SI{540}{\kilo\hertz} to \SI{1600}{\kilo\hertz} the frequency changes by a factor of 3, so $C$ by a factor of 9: that is why tuning capacitors have many interleaved plates.
\end{abcliste}
